\section{Introduction}\label{sec:intro}


In order to approach this problem, we revisit some early works of Mason \cite{MGlueing} and Las Vergnas \cite{LV81} on symmetric and tensor products of matroids. Both authors show that matroid powers do not necessarily exist. In the case of LasVergnas, it is shown that nontrivial tensor products with the V\'amos matroid never exist, and for Mason a specific point configuration is shown to lack large symmetric powers. Moreover, we can deduce from the nonexistence of matroid tensor powers the non-existence of second matroid symmetric powers; specifically, since the V\'amos matroid $V_8$ has no tensor powers, $V_8\oplus V_8$ has no second symmetric power.

We will follow the min-plus convention for the tropical semiring, which is the set $\T:= \R\cup \set{\infty}$ together with operations $\oplus$ and $\ttimes$ with $a\oplus b= \min(a,b)$ and $a\ttimes b = a+b$. Our approach in this paper naturally extends the work of Mason to the context of \textbf{valuated matroids}, which are subsemimodules $\sV\subset \T^n$ satisfying the following \textbf{vector elimination axiom}
    \begin{quote}
        For any vectors $G,H\in \sV$ and any $e\in [n]$ with $G_e=H_e$, there exists a vector $F\in \sV$ such that $F_e=\infty$, and $F_{e'}\geq\min{G_{e'},H_{e'}}$ for all $e'\neq e$ and with equality whenever $G_{e'}\neq H_{e'}$.
    \end{quote}
A \textbf{matroid} presented by its vectors is analogous to a valuated matroid for which the subsemimodule is generated by vectors $G$ such that $G_e\in \set{0,\infty}$ for all $e\in [n]$.

Throughout the paper we will make use of classical matroid theory and cryptomorphisms such as those for bases, circuits, spanning sets, flats, and hyperplanes as well as standard constructions such as taking minors and duals. We refer the reader to \cite{Oxley06} for a comprehensive background. We will define any relevant analogues for valuated matroids, such as the basis cryptomorphism and minors in \coloref{sec:LinTropIdeals}.

In this paper we will be primarily concerned with valuated matroid symmetric powers. Before providing this definition, we should discuss some useful notation. If $E$ is any set, the $d\th$ symmetric power of $E$, which we denote $\Sym_d(E)$, is the set of unordered $d$-tuples of elements of $E$. The motivation for defining $\Sym_d(E)$ comes from commutative algebra, where the set of degree $d$ monomials in $\set{x_1,\dots,x_n}$, denoted $\Mon_d(n)$, is precisely the set $\Sym_d(\set{x_1,\dots,x_n})$.

If we fix an element $\alpha \in \Sym_i(E)$, there is an injective ``multiplication" map
    \begin{align*}
        *_\alpha: \Sym_d(E)&\to \Sym_{d+i}(E)\\
        \beta&\mapsto (\alpha,\beta).
    \end{align*}
Since the coordinates of elements of $\Sym_{d+i}$ are unordered, we will write $(\alpha,\beta)$ as $\alpha\beta$, taking after commutative multiplication notation. 

We will often consider the subset of elements of $\Sym_{d+i}(E)$ that are divisible by some $\alpha\in \Sym_{i}(E)$, which we denote as

$$\Sym_{d+i}(E)|_{\alpha} := \set{e\in \Sym_{d+i}(E): e= \alpha x}.$$

Given a valuated matroid $\sV$ on $\Sym_{d+i}(E)$ we further extend this notation and write $\sV|_{\alpha}$ instead of $\sV|_{\Sym_{d+1}(E)|_{\alpha}},$ where the restriction of $\sV$ to the set $T\subset [n]$ is given by the subsemimodule 
$$S|_T:=\sV\cap \set{\pi_T(x)\in \T^T : x_i=\infty \text{ for all } i\notin T}.$$

\begin{definition}\label{def:introSymprod}
Fixing a matroid $\sV$ on a ground set $E$, a $\mathbf{d}^{\text{\textbf{th}}}$ \textbf{symmetric quasi power} of $\sV$ is a sequence of matroids $(\sV_1,\dots,\sV_d)$ with $\sV_1 = \sV$ and for all $1\leq i \leq d$
\begin{itemize}
    \item $E(\sV_i) = \Sym_i(E),$
    \item for all $\alpha\in \Sym_{d-i}(E)$ for which $\alpha$ is not divisible by a loop the multiplication map $*_\alpha$ given by $y\mapsto \alpha x$ induces the isomorphism
    \begin{equation} 
        \sV_i\cong \sV_d|_{\alpha}
    \end{equation}
    otherwise
    \begin{equation}
        \rk(\sV_d|_{\alpha})=0.
    \end{equation}
    
\end{itemize} 
Moreover, if $(\sV_1,\dots,\sV_d)$ is a symmetric quasi power of $\sV$ with

\begin{equation}
\rk(\sV_d)= \binom{\rk(\sV)+d-1}{d}
\end{equation}

then we say that $(\sV_1,\dots,\sV_d)$ is a $\mathbf{d}^{\text{\textbf{th}}}$ \textbf{symmetric power} of $\sV$.
\end{definition}

In \coloref{sec:MatProds} we will discuss relevant background on tensor and symmetric powers, establish geometric connections, and address how our definition differs from that of Mason in \cite{MGlueing}. We further prove the following theorem.

\begin{theorem}\label{thm:introThm0}
    The family of matroids with $d\th$ symmetric power is minor closed and has infinitely many forbidden minors.
\end{theorem}

\coloref{thm:introThm0} partly relies on the connection we make between symmetric powers and the geometry of their underlying matroids. To establish this connection we use the theory of tropical ideals.

\begin{definition}\label{def:tropicalIdeals}
    Let $I\subset \T[x_1,\dots,x_n]$ be a homogeneous ideal. Then $I$ is called a \textbf{tropical ideal} if $I_d:=I|_{\Mon_d(n)}$ is a valuated matroid in $\T^{\Mon_d(n)}$ for all $d\in \N$. 
\end{definition}

The most simple examples of tropical ideals are \textbf{realizable tropical ideals}, which are the tropicalizations of polynomial ideals $\trop(J)$, where $J$ is an ideal in $k[x_1,\dots,x_n]$ for some non-archimedean valued field $k$. We refer the reader to \cite{MR18} and \cite{MR20} for explicit examples, extensions to other polynomial semirings, and for a comprehensive background.

The underlying geometry of matroid symmetric powers and of tropical ideals is vital to establishing their connection. If $p\in \T[x_1,\dots,x_n]$ is a polynomial, its variety is given by
$$V(p)=\set{x\in \T^n: p \text{ attains its minimum at least twice at } x \text{ or } p(x)=\infty},$$
which is a balanced polyhedral complex in $\T^n$. The variety of a tropical ideal $I\subset \T[x_1,\dots,x_n]$ is then given by
$$V(I)=\bigcap_{p\in I}V(p),$$
which is also a balanced polyhedral complex, with weights defined as in \cite{MR20}. In this paper we refer to any balanced polyhedral complex in $\T^n$ as a \textbf{tropical variety}, and we distinguish those tropical varieties that arise as the variety of a tropical ideal by referring to them as \textbf{tropically realizable varieties}. 


To relate tropical geometry explicitly to valuated matroids, we consider that $\T^n\cong \T[x_1,\dots,x_n]_1$ under the map $e_i\mapsto x_i$. In particular, a valuated matroid $\sV\subset \T^n$ can be seen as a collection of polynomials in $\T[x_1,\dots,x_n].$ Explicitly, given a vector $v\in \sV$ we have the associated vector-polynomial
$$p_v = \bigoplus_{i\in [n]} v_ix_i\in \T[x_1,\dots,x_n].$$
The \textbf{tropical linear space} associated to $\sV$ is then given by
$$\trop(\sV) := \bigcap_{v\in \sV}V(p_v),$$
which is a balanced polyhedral complex in $\T^n$ with constant weights $1$.

\begin{theorem}\label{thm:introThm1}
Let $\sV$ be a matroid. The tropical linear space $\trop(\sV)$ is the variety of a tropical ideal if and only if a $d\th$ symmetric power of $\sV$ exists for every $d\in \N$.
\end{theorem}

If the variety $\trop(\sV)$ is of the form $V(I)$ for a tropical ideal $I$, as in \coloref{thm:introThm1}, then $\trop(\sV)$ is a tropically realizable variety by definition. We extend this terminology to the matroid $\sV$ itself and say that $\sV$ is a \textbf{tropically realizable matroid} whenever $\trop(\sV)$ is a tropically realizable variety. Tropically realizable varieties behave nicely under coordinate projections and stable intersections with hyperplanes \cite{MR20}, which generalize the standard matroid minor operations of deletion and contraction, respectively, to tropical ideals. We will further see that, in the case of tropically realizable matroids, these geometric operation extend naturally to matroid minor operations on symmetric powers.

\begin{theorem}\label{thm:introThm2}
Tropically realizable matroids are a minor closed class. In particular, if $\sV$ is a matroid on $E$ and $(\sV_1,\sV_2,\dots,\sV_d)$ is a $d\th$ symmetric power of $\sV$, then for any minor $\sV'$ of $\sV$, there is a minor $\sV'_i$ of each $\sV_i$ for which $(\sV',\sV'_2,\dots, \sV'_d)$ is a $d\th$ symmetric power of $\sV'$.
\end{theorem}

The geometric techniques used to prove \coloref{thm:introThm2} extend to other classes of matroids. We say that the matroid $\sV$ is \textbf{tropically set-realizable} if there is a tropical ideal $I$ for which the support of $V(I)$ is equal to that of $\trop(\sV)$, but the weights are not necessarily equal to $1$. Tropically set-realizable matroids were originally discussed in \cite{Yu17}, and here the the proof of \coloref{thm:introThm2} immediately implies that this set of matroids is minor closed. 

We also consider the class of tropically algebraic matroids. Given a (tropical) ideal $I$ in $\T[x_1,\dots,x_n]$, we define the independence complex of $I$ to be the collection of subsets of $[n]$ given by

$$\mathcal{I}(I):= \set{S\subset [n]: I\cap \T[x_i: i\in S]= \set{\infty}}.$$
The independence complex of $I$ is always an abstract simplicial complex, but it is not always a matroid, even if we restrict to tropical ideals $I$. We say that a matroid $M$ is \textbf{tropically algebraic} if there is a tropical ideal $I$ for which $\mathcal{I}(I)$ is the collection of independent sets of $M$. In \coloref{sec:LinTropIdeals} we will prove the set of tropically algebraic matroids is minor closed as well. 

Underlying these aforementioned proofs are several technical results that are quite useful, and in particular we'd highlight a special case of \coloref{prop:connectedCodim1}. If $\Sigma$ is a polyhedral complex then we can define its \textbf{facet-ridge} graph $G(\Sigma)$. This is the graph whose vertices are given by the facets (top-dimensional cells) of $\Sigma$ and where two vertices share an edge if there is a ridge (codimension-1 face) between the associated facets in $\Sigma$. The polyhedral complex $\Sigma$ is said to be connected through codimension one if $G(\Sigma)$ is a connected graph.

\begin{corollary}\label{thm:introThm3}
Tropical linear spaces are connected through codimension 1.
\end{corollary}

\coloref{thm:introThm3} is vital to the proof of \coloref{thm:introThm1}, and has implications well beyond the scope of this paper.

\section*{Acknowledgements}
I'd like to give some special thanks to Felipe Rinc\'on and Alex Fink for their valuable insight and guidance with this paper, as well as to Josephine Yu, Diane Maclagan, and Zach Walsh for helpful conversations that inspired multiple improvements to the results in this paper.
%%%%%%%%%%%%%%%%%%%%%%%%%%%%%%%%%%%%%%%%%%%%%%%%%%%%

%%%%%%%%%%%%%%%%%%%%%%%%%%%%%%%%%%%%%%%%%%%%%%%%%%%%

%%%%%%%%%%%%%%%%%%%%%%%%%%%%%%%%%%%%%%%%%%%%%%%%%%%%

    %%%%%%%%%%%%%%%%%%%%%%%%%%%%%%%%%%%%%%%%%%%%%%%%%%%%%%%%%%%
